\section*{Tipps \& Tricks}

\subsection*{Ordnung}
$o(g^m) = \frac{\kgV(m,o(g))}{|m|}$

\subsubsection*{In $\Z_n$}
$\bar m \in \Z_n$ Dann: $o(\bar m) = \frac{\kgV(m,n)}{m}$

\subsection*{Polynome}
\subsubsection*{Nullpolynom}
Hat ein Polynom mehr paarweise verschiedene Nullstellen als sein Grad,
ist es das Nullpolynom.

\subsection*{Permutationen}

$\displaystyle \sigma = \mat{1 & 2 & 3 & 4 & 5 \\ 3 & 4 & 5 & 2 & 1}$

\subsubsection*{Zykelschreibweise}
$\sigma = (1~3~5) \circ (2~4)$

\subsubsection*{Transpositionsschreibweise}
$\sigma = (1~3) \circ (3~5) \circ (2~4)$

Benachbarte Zahlen: \\
$\sigma = (1~2) \circ (2~3) \circ (2~1) \circ (3~4) \circ (4~5) \circ (4~3) \circ (2~3) \circ (3~4) \circ (3~2)$

\subsection*{Ringe Klassifikation}
$R$ Körper $\iff R[t]$ eukl.$\iff R[t]$ Hauptidealring

\subsection*{Einheiten in $\Z_n$}
$m$ ist Einheit, wenn $\ggT(m,n) = 1$.
